\documentclass[a4paper,12pt]{article}
\usepackage[utf8]{inputenc}
\usepackage[ngerman]{babel}
\usepackage{amsmath,amssymb}
\usepackage{geometry}\geometry{margin=2.5cm}
\usepackage{enumitem}
\usepackage{booktabs}
\usepackage{tikz}
\usetikzlibrary{arrows.meta, positioning, calc, decorations.pathreplacing}
\usepackage{pgfplots}\pgfplotsset{compat=1.18}
\usepackage{tcolorbox}
\tcbuselibrary{breakable}
\usepackage{xcolor}
\newtcolorbox{begriffe}[1][]{colback=yellow!10, colframe=yellow!60!black, fonttitle=\bfseries, title={Was bedeuten die Begriffe?}, breakable, #1}
\newtcolorbox{wastun}[1][]{colback=orange!5, colframe=orange!60!black, fonttitle=\bfseries, title={Was ist zu tun?}, breakable, #1}
\newcommand{\piv}[1]{{\color{red}\boxed{#1}}}
\newcommand{\grafik}{\textcolor{gray}{\small\textit{[Grafische Erkl\"arung -- nicht Teil der L\"osung]}}}
\newcommand{\erglerk}{\textcolor{gray}{\small\textit{[Erkl\"arung erg\"anzt]}}}
\newcommand{\EW}{Eigenwert}
\newcommand{\EV}{Eigenvektor}

\title{\"Ubungsblatt 11 -- L\"osungen \\ \large Linear Algebra}
\author{}
\date{}

\begin{document}
\maketitle

\begin{begriffe}
\begin{itemize}[leftmargin=*, nosep]
    \item \textbf{\EW{} $\lambda$ und \EV{} $v$}: Eine Zahl $\lambda$ und ein Vektor $v \neq 0$ mit $A v = \lambda v$. Das hei\ss t: $A$ \emph{streckt} den Vektor $v$ nur (um den Faktor $\lambda$), ohne seine Richtung zu \"andern.
    \item \textbf{Charakteristisches Polynom $\chi_A(t) = \det(A - t I)$}: Die Nullstellen dieses Polynoms sind genau die \EW{}e. (Begr\"undung: $A v = \lambda v$ hei\ss t $(A - \lambda I) v = 0$ mit $v \neq 0$, also ist $A - \lambda I$ \emph{nicht} invertierbar, also $\det(A - \lambda I) = 0$.)
    \item \textbf{Eigenraum $E_\lambda = \ker(A - \lambda I)$}: Alle \EV{}en zu $\lambda$ (plus der Nullvektor). Man findet ihn, indem man das LGS $(A - \lambda I) v = 0$ l\"ost.
    \item \textbf{Algebraische Vielfachheit AM($\lambda$)}: Wie oft $\lambda$ als Nullstelle im char. Polynom vorkommt (der Exponent von $(t - \lambda)$).
    \item \textbf{Geometrische Vielfachheit GM($\lambda$)}: $\dim E_\lambda$ $=$ Anzahl der Basisvektoren des Eigenraums $=$ Anzahl freier Variablen beim L\"osen von $(A - \lambda I)v = 0$.
    \item \textbf{Es gilt immer} $\;1 \le \text{GM}(\lambda) \le \text{AM}(\lambda)$. Sind alle GM $=$ AM, ist $A$ \emph{diagonalisierbar}.
\end{itemize}
\end{begriffe}

\grafik

\begin{center}
\begin{tikzpicture}[scale=1.0, >={Stealth[scale=1.1]}]
% axes
\draw[help lines, gray!25] (-0.3,-0.3) grid (4.2,3.2);
\draw[->, gray] (-0.3,0) -- (4.4,0) node[right, gray]{};
\draw[->, gray] (0,-0.3) -- (0,3.4) node[above, gray]{};
% eigen-direction: v and Av = lambda v on the SAME line
\draw[->, red!80, very thick] (0,0) -- (1.4,0.8) node[below right, red!80]{$v$};
\draw[->, red!50, very thick] (0,0) -- (2.8,1.6) node[above]{\,$Av = \lambda v$};
\draw[red!40, dashed] (1.4,0.8) -- (2.8,1.6);
% generic vector: w and Aw point in DIFFERENT directions
\draw[->, blue!80, very thick] (0,0) -- (0.7,2.2) node[above left, blue!80]{$w$};
\draw[->, blue!45, very thick] (0,0) -- (2.4,2.5) node[right, blue!70]{$Aw$};
\draw[blue!35, dashed] (0.7,2.2) -- (2.4,2.5);
\end{tikzpicture}
\end{center}

\textcolor{gray}{\small\textit{Die Grundidee: Multipliziert man einen Vektor mit $A$, wird er normalerweise gedreht \emph{und} gestreckt (blau: $w \to Aw$ zeigt woanders hin). Ein \textbf{\EV} (rot) ist eine besondere Richtung, die $A$ \emph{nur streckt} -- $Av$ liegt auf derselben Geraden wie $v$. Der Streckfaktor ist der \EW{} $\lambda$.}}

\bigskip
\hrule
\bigskip

%% ============================================================
\subsection*{Aufgabe 43}
%% ============================================================

\textbf{Angabe.} Bestimme f\"ur $A = \begin{pmatrix} 4 & 1 & 1 \\ 2 & 5 & 2 \\ -1 & -1 & 2 \end{pmatrix}$ und $B = \begin{pmatrix} 5 & 2 & 3 \\ 2 & 4 & 2 \\ -3 & -2 & -1 \end{pmatrix}$
das charakteristische Polynom, die \EW{}e mit algebraischer und geometrischer Vielfachheit sowie eine Basis jedes Eigenraums -- auf \textbf{zwei Wegen}.

\begin{wastun}
\begin{itemize}[leftmargin=*, nosep]
    \item \textbf{Weg 1:} $\chi_A(t) = \det(A - tI)$ mit der \textbf{Sarrus-Regel} ausmultiplizieren, dann mit dem \textbf{Satz \"uber ganzzahlige Nullstellen} faktorisieren (Kandidaten $=$ Teiler des konstanten Glieds).
    \item \textbf{Weg 2:} Durch \textbf{Zeilen-/Spaltenoperationen} fr\"uh einen Faktor $(\lambda - t)$ herausziehen -- das spart das Ausmultiplizieren.
    \item Danach pro \EW{} $\lambda$ das LGS $(A - \lambda I)v = 0$ l\"osen $\Rightarrow$ Basis des Eigenraums und GM.
\end{itemize}
\end{wastun}

%% --- 43 A ---
\subsubsection*{Matrix $A$}

\textbf{Weg 1: Sarrus $+$ ganzzahlige Nullstellen.}

\grafik

\begin{center}
\begin{tikzpicture}[scale=0.85, every node/.style={font=\small}]
% 3x3 matrix entries plus repeated first two columns
\foreach \x/\lab in {0/a,1/b,2/c,3/a,4/b} { \node at (\x,2) {$\lab$}; }
\foreach \x/\lab in {0/d,1/e,2/f,3/d,4/e} { \node at (\x,1) {$\lab$}; }
\foreach \x/\lab in {0/g,1/h,2/i,3/g,4/h} { \node at (\x,0) {$\lab$}; }
% separator after column 3
\draw[gray!40] (2.5,-0.5) -- (2.5,2.5);
% down diagonals (+) blue
\foreach \sx in {0,1,2} {
  \draw[->, blue!70, thick] (\sx-0.25,2.25) -- (\sx+1.75,-0.25);
}
% up diagonals (-) red
\foreach \sx in {0,1,2} {
  \draw[->, red!75, thick] (\sx-0.25,-0.25) -- (\sx+1.75,2.25);
}
\node[blue!70] at (1.0,-0.8) {$+\;(\text{nach unten})$};
\node[red!75] at (3.6,-0.8) {$-\;(\text{nach oben})$};
\end{tikzpicture}
\end{center}

\textcolor{gray}{\small\textit{Sarrus-Regel: Man schreibt die ersten zwei Spalten nochmal rechts daneben. Die drei Diagonalen \emph{nach unten} (blau) werden $+$ gez\"ahlt, die drei \emph{nach oben} (rot) $-$. Formel: $\det = aei + bfg + cdh - ceg - afh - bdi$.}}

\textbf{Schritt 1: $A - tI$ hinschreiben.} $tI$ ist die Diagonalmatrix mit lauter $t$ auf der Diagonale. Beim Abziehen \"andert sich also \emph{nur die Diagonale} von $A$:
\[
A - tI = \begin{pmatrix} 4 & 1 & 1 \\ 2 & 5 & 2 \\ -1 & -1 & 2 \end{pmatrix} - \begin{pmatrix} t & 0 & 0 \\ 0 & t & 0 \\ 0 & 0 & t \end{pmatrix} = \begin{pmatrix} 4-t & 1 & 1 \\ 2 & 5-t & 2 \\ -1 & -1 & 2-t \end{pmatrix}.
\]

\textbf{Schritt 2: Sarrus anwenden.}

\emph{Die Regel ganz allgemein.} F\"ur eine beliebige $3\times3$-Matrix mit Buchstaben gilt:
\[
\begin{pmatrix} a & b & c \\ d & e & f \\ g & h & i \end{pmatrix}
\;\longrightarrow\;
\det = aei + bfg + cdh \;-\; ceg - afh - bdi.
\]
(Die ersten drei Produkte $=$ Diagonalen nach unten $(+)$, die letzten drei $=$ Diagonalen nach oben $(-)$, siehe Grafik.)

\emph{Die Buchstaben aus unserer Matrix ablesen.} Aus $A - tI = \left(\begin{smallmatrix} 4-t & 1 & 1 \\ 2 & 5-t & 2 \\ -1 & -1 & 2-t \end{smallmatrix}\right)$ lesen wir Eintrag f\"ur Eintrag ab:
\[
\begin{array}{lll}
a = 4-t, & b = 1, & c = 1, \\
d = 2, & e = 5-t, & f = 2, \\
g = -1, & h = -1, & i = 2-t.
\end{array}
\]

\emph{Die sechs Produkte einzeln ausrechnen:}
\begin{align*}
aei &= (4-t)(5-t)(2-t), & ceg &= 1\cdot(5-t)\cdot(-1) = -(5-t), \\
bfg &= 1\cdot 2\cdot(-1) = -2, & afh &= (4-t)\cdot 2\cdot(-1) = -2(4-t), \\
cdh &= 1\cdot 2\cdot(-1) = -2, & bdi &= 1\cdot 2\cdot(2-t) = 2(2-t).
\end{align*}

\emph{In die Formel einsetzen} (erste drei mit $+$, letzte drei mit $-$):
\begin{align*}
\det(A - tI) &= aei + bfg + cdh - ceg - afh - bdi \\
&= (4-t)(5-t)(2-t) + (-2) + (-2) - \big(-(5-t)\big) - \big(-2(4-t)\big) - 2(2-t) \\
&= (4-t)(5-t)(2-t) - 2 - 2 + (5-t) + 2(4-t) - 2(2-t).
\end{align*}
Der Klammerterm ausmultipliziert: $(4-t)(5-t)(2-t) = -t^3 + 11t^2 - 38t + 40$. Die restlichen Terme: $-2 -2 +5 -t +8 -2t -4 +2t = 5 - t$. Zusammen:
\[
\chi_A(t) = -t^3 + 11t^2 - 38t + 40 + 5 - t = -t^3 + 11t^2 - 39t + 45.
\]

\erglerk\ \textcolor{gray}{\small Schnellcheck: konstantes Glied $\chi_A(0) = 45 = \det A$ \checkmark, und der $t^2$-Koeffizient $11 = \operatorname{Spur}(A) = 4+5+2$ \checkmark.}

\textbf{Schritt 3: Faktorisieren.} Mit $-1$ ausgeklammert: $\chi_A(t) = -(t^3 - 11t^2 + 39t - 45)$. Ganzzahlige Nullstellen sind Teiler von $45$ -- Ausprobieren zeigt: $t = 3$ und $t = 5$ ergeben $0$. Die dritte Nullstelle liefert die Spur ($3 + 5 + \lambda_3 = 11$): $\lambda_3 = 3$. Die Nullstellen sind also $\mathbf{3,\ 3,\ 5}$.

\emph{Von den Nullstellen zur Faktorform.} Jede Nullstelle $\lambda$ wird zu einem Faktor $(t - \lambda)$ -- eine doppelte Nullstelle entsprechend zum Quadrat:
\[
t^3 - 11t^2 + 39t - 45 = (t-3)\,\underbrace{(t-3)}_{\text{die } 3 \text{ doppelt}}(t-5) = (t-3)^2(t-5).
\]
Das $-1$ von vorne ziehen wir in den letzten Faktor hinein, und wegen $(t-3)^2 = (3-t)^2$ sowie $-(t-5) = (5-t)$ wird daraus:
\[
\boxed{\chi_A(t) = -(t-3)^2(t-5) = (3-t)^2(5-t).}
\]

\emph{Exponent $=$ algebraische Vielfachheit.} Die AM ist einfach, wie oft die Nullstelle vorkommt $=$ der Exponent des Faktors: $\lambda = 3$ steht im Quadrat $\Rightarrow$ AM $2$; $\lambda = 5$ steht einfach $\Rightarrow$ AM $1$.

\bigskip

\textbf{Weg 2: Faktor fr\"uh herausziehen (Spaltentrick).}

\grafik

\begin{center}
\begin{tikzpicture}[scale=0.85, every node/.style={font=\small}]
\node[anchor=west] (M) at (0,0) {$\begin{pmatrix} 4{-}t & 1 & 1 \\ 2 & 5{-}t & 2 \\ -1 & -1 & 2{-}t \end{pmatrix}$};
\draw[->, orange!80!black, thick] (3.6,0) -- (5.3,0);
\node[orange!80!black] at (4.45,0.45) {\footnotesize $Z_1\!+\!Z_2\!+\!Z_3$};
\node[anchor=west] (R) at (5.6,0) {$\begin{pmatrix} 5{-}t & 5{-}t & 5{-}t \\ 2 & 5{-}t & 2 \\ -1 & -1 & 2{-}t \end{pmatrix}$};
\end{tikzpicture}
\end{center}

\textcolor{gray}{\small\textit{Beobachtung: In $A$ summiert jede Spalte zu $5$ (z.\,B. $4+2-1 = 5$). Also summiert in $A - tI$ jede Spalte zu $5 - t$. Wenn man alle Zeilen aufaddiert ($Z_1 \to Z_1 + Z_2 + Z_3$), entsteht oben eine Zeile aus lauter $(5-t)$ -- die kann man als Faktor herausziehen.}}

\textbf{Schritt 1: $Z_1 \to Z_1 + Z_2 + Z_3$} (\"andert die Determinante nicht). Wir addieren die drei Zeilen von $A - tI$ und schreiben das Ergebnis in die erste Zeile -- spaltenweise:
\begin{align*}
\text{Spalte 1:}\;\; (4-t) + 2 + (-1) &= 5 - t, \\
\text{Spalte 2:}\;\; 1 + (5-t) + (-1) &= 5 - t, \\
\text{Spalte 3:}\;\; 1 + 2 + (2-t) &= 5 - t.
\end{align*}
Die neue erste Zeile ist also $(5-t,\, 5-t,\, 5-t)$, die anderen beiden bleiben. Jetzt $(5-t)$ ausklammern:
\[
\det(A - tI) = \det\!\begin{pmatrix} 5-t & 5-t & 5-t \\ 2 & 5-t & 2 \\ -1 & -1 & 2-t \end{pmatrix}
= (5-t)\,\det\!\begin{pmatrix} 1 & 1 & 1 \\ 2 & 5-t & 2 \\ -1 & -1 & 2-t \end{pmatrix}.
\]

\textbf{Schritt 2: Spalten $C_2 \to C_2 - C_1$ und $C_3 \to C_3 - C_1$} (Spalte 1 von Spalte 2 und 3 abziehen, um oben rechts Nullen zu bekommen). Eintragweise:
\begin{align*}
C_2 - C_1:\;\; & 1-1 = 0,\;\; (5-t)-2 = 3-t,\;\; -1-(-1) = 0 \;\Rightarrow\; \text{Spalte } (0,\, 3-t,\, 0), \\
C_3 - C_1:\;\; & 1-1 = 0,\;\; 2-2 = 0,\;\; (2-t)-(-1) = 3-t \;\Rightarrow\; \text{Spalte } (0,\, 0,\, 3-t).
\end{align*}
\[
= (5-t)\,\det\!\begin{pmatrix} 1 & 0 & 0 \\ 2 & 3-t & 0 \\ -1 & 0 & 3-t \end{pmatrix}
= (5-t)\,(3-t)(3-t) = (5-t)(3-t)^2.
\]

\erglerk\ \textcolor{gray}{\small Die letzte Matrix ist (untere) Dreiecksform $\Rightarrow$ Determinante $=$ Produkt der Diagonale $1\cdot(3-t)\cdot(3-t)$.}

Gleiches Ergebnis wie Weg 1: $\chi_A(t) = (5-t)(3-t)^2$, \EW{}e $\lambda = 5$ (AM $1$), $\lambda = 3$ (AM $2$).

\medskip
\textbf{Eigenr\"aume.}

\textit{$\lambda = 5$:} L\"ose $(A - 5I)v = 0$. Zuerst $A - 5I$ (nur die Diagonale wird um $5$ kleiner):
\[
A - 5I = \begin{pmatrix} 4 & 1 & 1 \\ 2 & 5 & 2 \\ -1 & -1 & 2 \end{pmatrix} - \begin{pmatrix} 5 & 0 & 0 \\ 0 & 5 & 0 \\ 0 & 0 & 5 \end{pmatrix} = \begin{pmatrix} -1 & 1 & 1 \\ 2 & 0 & 2 \\ -1 & -1 & -3 \end{pmatrix}.
\]
Dann Gau\ss:
\[
\begin{pmatrix} -1 & 1 & 1 \\ 2 & 0 & 2 \\ -1 & -1 & -3 \end{pmatrix}
\xrightarrow[Z_3 \to Z_3 - Z_1]{Z_2 \to Z_2 + 2Z_1}
\begin{pmatrix} \piv{-1} & 1 & 1 \\ 0 & \piv{2} & 4 \\ 0 & -2 & -4 \end{pmatrix}
\xrightarrow{Z_3 \to Z_3 + Z_2}
\begin{pmatrix} \piv{-1} & 1 & 1 \\ 0 & \piv{2} & 4 \\ 0 & 0 & 0 \end{pmatrix}.
\]
Jede Zeile der Stufenform ist eine Gleichung f\"ur $(x,y,z)$ -- Eintr\"age mal $x, y, z$, gleich $0$: \; $(0,2,4) \Rightarrow 2y + 4z = 0$, \quad $(-1,1,1) \Rightarrow -x + y + z = 0$. Freie Variable ist $z$; aufl\"osen: $2y + 4z = 0 \Rightarrow y = -2z$, dann $-x + y + z = 0 \Rightarrow x = y + z = -z$. Mit $z = 1$:
\[
E_5 = \operatorname{span}\{(-1, -2, 1)\}, \qquad \text{GM}(5) = 1.
\]

\textit{$\lambda = 3$:} Zuerst $A - 3I$:
\[
A - 3I = \begin{pmatrix} 4 & 1 & 1 \\ 2 & 5 & 2 \\ -1 & -1 & 2 \end{pmatrix} - \begin{pmatrix} 3 & 0 & 0 \\ 0 & 3 & 0 \\ 0 & 0 & 3 \end{pmatrix} = \begin{pmatrix} 1 & 1 & 1 \\ 2 & 2 & 2 \\ -1 & -1 & -1 \end{pmatrix}.
\]
Dann Gau\ss:
\[
\begin{pmatrix} 1 & 1 & 1 \\ 2 & 2 & 2 \\ -1 & -1 & -1 \end{pmatrix}
\xrightarrow[Z_3 \to Z_3 + Z_1]{Z_2 \to Z_2 - 2Z_1}
\begin{pmatrix} \piv{1} & 1 & 1 \\ 0 & 0 & 0 \\ 0 & 0 & 0 \end{pmatrix}.
\]
Die einzige nichttriviale Zeile $(1,1,1)$ bedeutet $x + y + z = 0$ (Eintr\"age mal $x, y, z$). Freie Variablen $y, z$, also $x = -y - z$. Mit $(y,z) = (1,0)$ und $(0,1)$:
\[
E_3 = \operatorname{span}\{(-1, 1, 0),\, (-1, 0, 1)\}, \qquad \text{GM}(3) = 2.
\]

\[
\fbox{\parbox{0.93\textwidth}{\centering
$\chi_A(t) = (5-t)(3-t)^2$. \\[3pt]
$\lambda = 5$: AM $1$, GM $1$, Basis $\{(-1,-2,1)\}$. \quad
$\lambda = 3$: AM $2$, GM $2$, Basis $\{(-1,1,0),(-1,0,1)\}$. \\[3pt]
Alle GM $=$ AM $\Rightarrow$ $A$ ist diagonalisierbar.
}}
\]

%% --- 43 B ---
\subsubsection*{Matrix $B$}

\textbf{Weg 1: Sarrus.} Zuerst $B - tI$ (wieder nur die Diagonale):
\[
B - tI = \begin{pmatrix} 5 & 2 & 3 \\ 2 & 4 & 2 \\ -3 & -2 & -1 \end{pmatrix} - \begin{pmatrix} t & 0 & 0 \\ 0 & t & 0 \\ 0 & 0 & t \end{pmatrix} = \begin{pmatrix} 5-t & 2 & 3 \\ 2 & 4-t & 2 \\ -3 & -2 & -1-t \end{pmatrix}.
\]
\emph{Buchstaben ablesen} (gleiche Formel $\det = aei + bfg + cdh - ceg - afh - bdi$ wie bei $A$):
\[
\begin{array}{lll}
a = 5-t, & b = 2, & c = 3, \\
d = 2, & e = 4-t, & f = 2, \\
g = -3, & h = -2, & i = -1-t.
\end{array}
\]

\emph{Die sechs Produkte:}
\begin{align*}
aei &= (5-t)(4-t)(-1-t), & ceg &= 3\cdot(4-t)\cdot(-3) = -9(4-t), \\
bfg &= 2\cdot 2\cdot(-3) = -12, & afh &= (5-t)\cdot 2\cdot(-2) = -4(5-t), \\
cdh &= 3\cdot 2\cdot(-2) = -12, & bdi &= 2\cdot 2\cdot(-1-t) = 4(-1-t).
\end{align*}

\emph{In die Formel einsetzen:}
\begin{align*}
\chi_B(t) &= aei + bfg + cdh - ceg - afh - bdi \\
&= (5-t)(4-t)(-1-t) + (-12) + (-12) - \big(-9(4-t)\big) - \big(-4(5-t)\big) - 4(-1-t) \\
&= (5-t)(4-t)(-1-t) - 12 - 12 + 9(4-t) + 4(5-t) - 4(-1-t) \\
&= (5-t)(4-t)(-1-t) + 36 - 9t = -t^3 + 8t^2 - 20t + 16.
\end{align*}

\erglerk\ \textcolor{gray}{\small Check: $\chi_B(0) = 16 = \det B$ \checkmark, $t^2$-Koeffizient $8 = \operatorname{Spur}(B) = 5+4-1$ \checkmark.}

\textbf{Faktorisieren.} $\chi_B(t) = -(t^3 - 8t^2 + 20t - 16)$. Ganzzahlige Nullstellen k\"onnen nur Teiler von $16$ sein ($\pm1, \pm2, \pm4, \dots$). Man probiert sie durch -- \emph{testen} hei\ss t: f\"ur $t$ einsetzen und schauen, ob $0$ herauskommt:
\[
t = 2:\; 8 - 32 + 40 - 16 = 0\ \checkmark, \qquad t = 4:\; 64 - 128 + 80 - 16 = 0\ \checkmark.
\]
$2$ und $4$ ergeben $0$, sind also Nullstellen. Die \textbf{dritte} (noch unbekannt, hei\ss e $\lambda_3$) ist geschenkt: die drei Nullstellen summieren sich zu $8$ (der Spur), und $2 + 4 + \lambda_3 = 8$ gibt $\lambda_3 = 2$. Nullstellen also $\mathbf{2,\ 2,\ 4}$. Jede Nullstelle wird zu einem Faktor, die doppelte $2$ zum Quadrat:
\[
\boxed{\chi_B(t) = -(t-2)^2(t-4) = (2-t)^2(4-t).}
\]
Exponent $=$ AM: $\lambda = 2$ im Quadrat $\Rightarrow$ AM $2$; $\lambda = 4$ einfach $\Rightarrow$ AM $1$.

\textbf{Weg 2: Faktor fr\"uh herausziehen (Spaltentrick).} Auch in $B$ summiert jede Spalte zu $4$ ($5+2-3=4$, $2+4-2=4$, $3+2-1=4$), in $B - tI$ also zu $4 - t$.

\textbf{Schritt 1: $Z_1 \to Z_1 + Z_2 + Z_3$} (\"andert die Determinante nicht). Die drei Zeilen von $B - tI$ aufaddiert, spaltenweise:
\begin{align*}
\text{Spalte 1:}\;\; (5-t) + 2 + (-3) &= 4 - t, \\
\text{Spalte 2:}\;\; 2 + (4-t) + (-2) &= 4 - t, \\
\text{Spalte 3:}\;\; 3 + 2 + (-1-t) &= 4 - t.
\end{align*}
Neue erste Zeile $(4-t,\, 4-t,\, 4-t)$, die anderen bleiben. Jetzt $(4-t)$ ausklammern:
\[
\det(B - tI) = \det\!\begin{pmatrix} 4-t & 4-t & 4-t \\ 2 & 4-t & 2 \\ -3 & -2 & -1-t \end{pmatrix}
= (4-t)\,\det\!\begin{pmatrix} 1 & 1 & 1 \\ 2 & 4-t & 2 \\ -3 & -2 & -1-t \end{pmatrix}.
\]

\textbf{Schritt 2: Spalten $C_2 \to C_2 - C_1$ und $C_3 \to C_3 - C_1$} (Spalte 1 von Spalte 2 und 3 abziehen). Eintragweise:
\begin{align*}
C_2 - C_1:\;\; & 1-1 = 0,\;\; (4-t)-2 = 2-t,\;\; -2-(-3) = 1 \;\Rightarrow\; \text{Spalte } (0,\, 2-t,\, 1), \\
C_3 - C_1:\;\; & 1-1 = 0,\;\; 2-2 = 0,\;\; (-1-t)-(-3) = 2-t \;\Rightarrow\; \text{Spalte } (0,\, 0,\, 2-t).
\end{align*}
\[
= (4-t)\,\det\!\begin{pmatrix} 1 & 0 & 0 \\ 2 & 2-t & 0 \\ -3 & 1 & 2-t \end{pmatrix}
= (4-t)\,(2-t)(2-t) = (4-t)(2-t)^2.
\]

\erglerk\ \textcolor{gray}{\small Die letzte Matrix ist (untere) Dreiecksform $\Rightarrow$ Determinante $=$ Produkt der Diagonale $1\cdot(2-t)\cdot(2-t)$.}

Gleiches Ergebnis wie Weg 1: $\chi_B(t) = (4-t)(2-t)^2$, \EW{}e $\lambda = 2$ (AM $2$), $\lambda = 4$ (AM $1$).

\medskip
\textbf{Eigenr\"aume.}

\textit{$\lambda = 4$:} Zuerst $B - 4I$:
\[
B - 4I = \begin{pmatrix} 5 & 2 & 3 \\ 2 & 4 & 2 \\ -3 & -2 & -1 \end{pmatrix} - \begin{pmatrix} 4 & 0 & 0 \\ 0 & 4 & 0 \\ 0 & 0 & 4 \end{pmatrix} = \begin{pmatrix} 1 & 2 & 3 \\ 2 & 0 & 2 \\ -3 & -2 & -5 \end{pmatrix}.
\]
Dann Gau\ss:
\[
\begin{pmatrix} 1 & 2 & 3 \\ 2 & 0 & 2 \\ -3 & -2 & -5 \end{pmatrix}
\xrightarrow[Z_3 \to Z_3 + 3Z_1]{Z_2 \to Z_2 - 2Z_1}
\begin{pmatrix} \piv{1} & 2 & 3 \\ 0 & \piv{-4} & -4 \\ 0 & 4 & 4 \end{pmatrix}
\xrightarrow{Z_3 \to Z_3 + Z_2}
\begin{pmatrix} \piv{1} & 2 & 3 \\ 0 & \piv{-4} & -4 \\ 0 & 0 & 0 \end{pmatrix}.
\]
Zeilen als Gleichungen: $(0,-4,-4) \Rightarrow -4y - 4z = 0 \Rightarrow y = -z$; \; $(1,2,3) \Rightarrow x + 2y + 3z = 0 \Rightarrow x = -2y - 3z = -z$. Mit $z=1$:
\[
E_4 = \operatorname{span}\{(-1,-1,1)\}, \qquad \text{GM}(4) = 1.
\]

\textit{$\lambda = 2$:} Zuerst $B - 2I$:
\[
B - 2I = \begin{pmatrix} 5 & 2 & 3 \\ 2 & 4 & 2 \\ -3 & -2 & -1 \end{pmatrix} - \begin{pmatrix} 2 & 0 & 0 \\ 0 & 2 & 0 \\ 0 & 0 & 2 \end{pmatrix} = \begin{pmatrix} 3 & 2 & 3 \\ 2 & 2 & 2 \\ -3 & -2 & -3 \end{pmatrix}.
\]
Dann Gau\ss:
\[
\begin{pmatrix} 3 & 2 & 3 \\ 2 & 2 & 2 \\ -3 & -2 & -3 \end{pmatrix}
\xrightarrow{Z_2 \to \tfrac12 Z_2}
\begin{pmatrix} 3 & 2 & 3 \\ 1 & 1 & 1 \\ -3 & -2 & -3 \end{pmatrix}
\xrightarrow[Z_3 \to Z_3 + 3Z_2']{Z_1 \to Z_1 - 3Z_2'}
\begin{pmatrix} 0 & -1 & 0 \\ \piv{1} & 1 & 1 \\ 0 & 1 & 0 \end{pmatrix}.
\]
Zeilen als Gleichungen: $(0,-1,0) \Rightarrow -y = 0$, also $y = 0$; \; $(1,1,1) \Rightarrow x + y + z = 0 \Rightarrow x = -z$. Freie Variable nur $z$. Mit $z=1$:
\[
E_2 = \operatorname{span}\{(-1, 0, 1)\}, \qquad \text{GM}(2) = 1.
\]

\[
\fbox{\parbox{0.93\textwidth}{\centering
$\chi_B(t) = (2-t)^2(4-t)$. \\[3pt]
$\lambda = 2$: AM $2$, GM $1$, Basis $\{(-1,0,1)\}$. \quad
$\lambda = 4$: AM $1$, GM $1$, Basis $\{(-1,-1,1)\}$. \\[3pt]
Hier ist GM$(2) = 1 < {}$AM$(2) = 2 \Rightarrow$ $B$ ist \emph{nicht} diagonalisierbar.
}}
\]

\bigskip
\hrule
\bigskip

%% ============================================================
\subsection*{Aufgabe 44}
%% ============================================================

\textbf{Angabe.} Bestimme \EW{}e (mit AM, GM) und Basen der Eigenr\"aume von
$A = \begin{pmatrix} 1 & 0 & 5 \\ 0 & 6 & 0 \\ 1 & 0 & 5 \end{pmatrix}$ und
$B = \begin{pmatrix} 5 & 5 & 5 & 5 \\ 7 & 7 & 7 & 7 \\ 11 & 11 & 11 & 11 \\ -23 & -23 & -23 & -23 \end{pmatrix}$.

\begin{wastun}
\begin{itemize}[leftmargin=*, nosep]
    \item \textbf{Kein charakteristisches Polynom!} Stattdessen Eigenschaften nutzen: $\lambda = 0$ genau dann, wenn $A$ nicht invertierbar ist; Summe der \EW{}e $= \operatorname{Spur}$; Produkt $= \det$; $\text{GM} \le \text{AM}$; offensichtliche \EV{}en ablesen.
\end{itemize}
\end{wastun}

%% --- 44 A ---
\subsubsection*{Matrix $A$}

\grafik

\begin{center}
\begin{tikzpicture}[scale=0.9, every node/.style={font=\small}]
\node at (0,0) {$A = \begin{pmatrix} 1 & 0 & 5 \\ 0 & 6 & 0 \\ 1 & 0 & 5 \end{pmatrix}$};
% highlight rows 1 and 3 equal
\draw[red!70, rounded corners] (-1.15,0.5) rectangle (1.15,0.95);
\draw[red!70, rounded corners] (-1.15,-0.95) rectangle (1.15,-0.5);
\node[red!70, right] at (1.5,0.72) {\footnotesize Zeile 1};
\node[red!70, right] at (1.5,-0.72) {\footnotesize $=$ Zeile 3 $\Rightarrow \det = 0 \Rightarrow \lambda = 0$};
\node[blue!70, right] at (1.5,0.0) {\footnotesize $A e_2 = (0,6,0) = 6\,e_2 \Rightarrow \lambda = 6$};
\end{tikzpicture}
\end{center}

\textcolor{gray}{\small\textit{Zwei Sachen sieht man sofort: (1) Zeile $1 =$ Zeile $3$, also ist $A$ singul\"ar $\Rightarrow$ $0$ ist \EW{}. (2) Die mittlere Spalte ist $(0,6,0)$, d.\,h. $A$ schickt $e_2 = (0,1,0)$ auf $6\,e_2$ $\Rightarrow$ $6$ ist \EW{} mit \EV{} $e_2$.}}

\textbf{Schritt 1: $\lambda = 0$.} Zeile $1 =$ Zeile $3$, und zwei gleiche Zeilen geben $\det A = 0$. Also ist $A$ nicht invertierbar, d.\,h. es gibt $v \neq 0$ mit $Av = 0 = 0\cdot v$ -- $\lambda = 0$ ist \EW{}, und der Eigenraum ist $\ker A$. Die Zeilen von $Av = 0$ als Gleichungen:
\[
\text{Zeile 2}\ (0,6,0): \ 6y = 0 \Rightarrow y = 0; \qquad \text{Zeile 1}\ (1,0,5): \ x + 5z = 0 \Rightarrow x = -5z.
\]
Frei ist $z$; mit $z = 1$:
\[
E_0 = \operatorname{span}\{(-5, 0, 1)\}, \qquad \text{GM}(0) = 1.
\]

\textbf{Schritt 2: $\lambda = 6$.} Die mittlere Spalte ist $(0,6,0)$, d.\,h. $A e_2 = 6 e_2$ -- also ist $6$ ein \EW{} mit \EV{} $e_2$. Den ganzen Eigenraum gibt $\ker(A - 6I)$:
\[
A - 6I = \begin{pmatrix} 1 & 0 & 5 \\ 0 & 6 & 0 \\ 1 & 0 & 5 \end{pmatrix} - \begin{pmatrix} 6 & 0 & 0 \\ 0 & 6 & 0 \\ 0 & 0 & 6 \end{pmatrix} = \begin{pmatrix} -5 & 0 & 5 \\ 0 & 0 & 0 \\ 1 & 0 & -1 \end{pmatrix}.
\]
Zeilen als Gleichungen: $(-5,0,5) \Rightarrow -5x + 5z = 0 \Rightarrow x = z$; $(1,0,-1) \Rightarrow x - z = 0$ (dasselbe). Auf $y$ keine Bedingung, also sind $y$ und $z$ frei. Mit $(y,z) = (1,0)$ und $(0,1)$:
\[
E_6 = \operatorname{span}\{(0,1,0),\, (1,0,1)\}, \qquad \text{GM}(6) = 2.
\]

\textbf{Schritt 3: Vielfachheiten.} Eine $3\times3$-Matrix hat insgesamt $3$ \EW{}e (mit Vielfachheit), die AMs summieren sich also zu $3$. Wegen GM $\le$ AM ist $\text{AM}(6) \ge \text{GM}(6) = 2$, und $0$ ist \EW{}, also $\text{AM}(0) \ge 1$. Schon $2 + 1 = 3$ -- damit ist alles vergeben: $\text{AM}(6) = 2$, $\text{AM}(0) = 1$.
\erglerk\ \textcolor{gray}{\small Probe mit der Spur: $0\cdot 1 + 6 \cdot 2 = 12 = \operatorname{Spur}(A) = 1+6+5$ \checkmark.}

\[
\fbox{\parbox{0.93\textwidth}{\centering
$\lambda = 0$: AM $1$, GM $1$, Basis $\{(-5,0,1)\}$. \quad
$\lambda = 6$: AM $2$, GM $2$, Basis $\{(0,1,0),(1,0,1)\}$. \\[3pt]
Alle GM $=$ AM $\Rightarrow$ $A$ ist diagonalisierbar.
}}
\]

%% --- 44 B ---
\subsubsection*{Matrix $B$}

\grafik

\begin{center}
\begin{tikzpicture}[scale=0.9, every node/.style={font=\small}]
\node[anchor=west] (Bm) at (0,0) {$B = \begin{pmatrix} 5 & 5 & 5 & 5 \\ 7 & 7 & 7 & 7 \\ 11 & 11 & 11 & 11 \\ -23 & -23 & -23 & -23 \end{pmatrix}$};
\draw[decorate, decoration={brace, amplitude=5pt}, orange!80!black] ([xshift=5pt, yshift=1.45cm]Bm.east) -- ([xshift=5pt, yshift=-1.45cm]Bm.east);
\node[orange!80!black, right] at ([xshift=13pt, yshift=0.6cm]Bm.east) {\footnotesize jede Zeile $= c_i \cdot (1,1,1,1)$};
\node[orange!80!black, right] at ([xshift=13pt, yshift=0.0cm]Bm.east) {\footnotesize $c = (5,7,11,-23)$};
\node[orange!80!black, right] at ([xshift=13pt, yshift=-0.6cm]Bm.east) {\footnotesize $\Rightarrow$ Rang $1$};
\end{tikzpicture}
\end{center}

\textcolor{gray}{\small\textit{$B$ hat \"uberall identische Spalten -- jede Zeile ist ein Vielfaches von $(1,1,1,1)$. So eine Matrix hat \textbf{Rang $1$}: alle Zeilen liegen ``auf einer Linie''. Damit ist der Kern riesig (Dimension $4 - 1 = 3$) und $0$ ist \EW{}.}}

\textbf{Schritt 1: $\lambda = 0$.} Jede Zeile von $B$ ist ein Vielfaches von $(1,1,1,1)$ -- es gibt also nur \emph{eine} unabh\"angige Zeile, $\operatorname{Rang}(B) = 1$. Mit der Dimensionsformel $\dim\ker(B) = 4 - \operatorname{Rang} = 4 - 1 = 3$. Ein Kern gr\"o\ss er als $\{0\}$ hei\ss t: es gibt $v \neq 0$ mit $Bv = 0$, also ist $0$ \EW{} mit $\text{GM}(0) = \dim\ker = 3$.

Der Kern: jede Zeile von $Bv = 0$ ist (z.\,B. Zeile 1) $5a + 5b + 5c + 5d = 5(a+b+c+d) = 0$, also $a + b + c + d = 0$. Eine Gleichung, vier Unbekannte $\Rightarrow$ drei freie Variablen ($b, c, d$; \, $a = -b-c-d$):
\[
E_0 = \operatorname{span}\{(-1,1,0,0),\, (-1,0,1,0),\, (-1,0,0,1)\}, \qquad \text{GM}(0) = 3.
\]

\textbf{Schritt 2: $\text{AM}(0) = 4$.} $B$ ist $4\times4$, hat also $4$ \EW{}e (mit Vielfachheit). Drei davon sind $0$ (denn $\text{AM}(0) \ge \text{GM}(0) = 3$). Der vierte hei\ss e $\mu$. Die Summe aller \EW{}e ist die Spur: $5 + 7 + 11 - 23 = 0$, also $0 + 0 + 0 + \mu = 0 \Rightarrow \mu = 0$. Damit ist $0$ der einzige \EW{}:
\[
\text{AM}(0) = 4.
\]

\erglerk\ \textcolor{gray}{\small Anders gesagt: Eine Rang-$1$-Matrix $c\,r^\top$ hat nur den \EW{} $r^\top c = \operatorname{Spur}$ (einfach) und sonst $0$. Hier ist $\operatorname{Spur} = 0$, also sind \emph{alle} \EW{}e $0$ (die Matrix ist nilpotent, $B^2 = 0$).}

\[
\fbox{\parbox{0.93\textwidth}{\centering
$\lambda = 0$: AM $4$, GM $3$, Basis $\{(-1,1,0,0),(-1,0,1,0),(-1,0,0,1)\}$. \\[3pt]
GM$(0) = 3 < {}$AM$(0) = 4 \Rightarrow$ $B$ ist \emph{nicht} diagonalisierbar.
}}
\]

\bigskip
\hrule
\bigskip

%% ============================================================
\subsection*{Aufgabe 45}
%% ============================================================

\textbf{Angabe.} $A, B \in \mathbb{F}^{n\times n}$. Zeige die folgenden Aussagen.

\begin{wastun}
\begin{itemize}[leftmargin=*, nosep]
    \item Bei jeder Aussage starten wir mit ``$\lambda$ ist \EW{} von $A$'' $=$ ``es gibt $v \neq 0$ mit $Av = \lambda v$'' und rechnen nach, dass derselbe $v$ auch \EV{} der neuen Matrix ist.
\end{itemize}
\end{wastun}

\textbf{(a) Ist $\lambda$ \EW{} von $A$ und $c \in \mathbb{F}$, so ist $\lambda + c$ \EW{} von $A + cI$.}

\textit{Beweis.} Sei $v \neq 0$ mit $Av = \lambda v$. Dann
\[
(A + cI)v = Av + c\,v = \lambda v + c v = (\lambda + c)\,v.
\]
Da $v \neq 0$, ist $\lambda + c$ \EW{} von $A + cI$ (mit demselben \EV{} $v$). $\qquad\square$

\textbf{(b) Ist $\lambda$ \EW{} von $A$, so ist $\lambda^2$ \EW{} von $A^2$.}

\textit{Beweis.} Sei $v \neq 0$ mit $Av = \lambda v$. Dann
\[
A^2 v = A(Av) = A(\lambda v) = \lambda (Av) = \lambda (\lambda v) = \lambda^2 v. \qquad\square
\]

\textbf{(c) Ist $A$ invertierbar und $\lambda$ \EW{} von $A$, so ist $\lambda^{-1}$ \EW{} von $A^{-1}$.}

\textit{Beweis.} Da $A$ invertierbar ist, ist $0$ kein \EW{} (sonst g\"abe es $v \neq 0$ mit $Av = 0$, also w\"are $A$ nicht injektiv). Also $\lambda \neq 0$. Aus $Av = \lambda v$ folgt durch Anwenden von $A^{-1}$:
\[
v = A^{-1}(\lambda v) = \lambda\, A^{-1} v \;\Longrightarrow\; A^{-1} v = \tfrac{1}{\lambda}\, v.
\]
Da $v \neq 0$, ist $\lambda^{-1}$ \EW{} von $A^{-1}$. $\qquad\square$

\textbf{(d) Ist $v$ \EV{} von $A$ \emph{und} von $B$, so ist $v$ auch \EV{} von $AB$.}

\textit{Beweis.} Sei $v \neq 0$ mit $Av = \lambda v$ und $Bv = \mu v$. Dann
\[
(AB)v = A(Bv) = A(\mu v) = \mu (Av) = \mu (\lambda v) = (\lambda\mu)\,v.
\]
Da $v \neq 0$, ist $v$ \EV{} von $AB$ zum \EW{} $\lambda\mu$. $\qquad\square$

\bigskip
\hrule
\bigskip

%% ============================================================
\subsection*{\"Uberblick: Ergebnisse 43 \& 44}
%% ============================================================

\begin{center}
\begin{tabular}{@{}llll@{}}
\toprule
Matrix & \EW{} ($\lambda$) & AM, GM & Basis des Eigenraums \\
\midrule
$43\,A$ & $3$ & $2,\,2$ & $(-1,1,0),\,(-1,0,1)$ \\
        & $5$ & $1,\,1$ & $(-1,-2,1)$ \\[2pt]
$43\,B$ & $2$ & $2,\,1$ & $(-1,0,1)$ \\
        & $4$ & $1,\,1$ & $(-1,-1,1)$ \\[2pt]
$44\,A$ & $0$ & $1,\,1$ & $(-5,0,1)$ \\
        & $6$ & $2,\,2$ & $(0,1,0),\,(1,0,1)$ \\[2pt]
$44\,B$ & $0$ & $4,\,3$ & $(-1,1,0,0),\,(-1,0,1,0),\,(-1,0,0,1)$ \\
\bottomrule
\end{tabular}
\end{center}

\textit{Diagonalisierbar} sind $43\,A$ und $44\,A$ (alle GM $=$ AM). \emph{Nicht} diagonalisierbar sind $43\,B$ und $44\,B$ (je ein GM $<$ AM).

\end{document}
